\Needspace{7\baselineskip}
\Article{316}
Following these principles, we have constructed a table of the necessary periods for all denominators of the form \(p^\mu\) below \(1000\). We will publish it in full, or even extended further, when an occasion arises. Here Table III, carried only as far as \(100\), may suffice as a sample; it hardly needs explanation. For denominators for which \(10\) is a primitive root, it gives the periods of fractions with numerator \(1\), namely for \(7,17,19,23,29,47,59,61,97\). For the others it gives \(f\) periods corresponding to the numerators \(1,r,rr,\ldots r^{f-1}\), distinguished by the attached numbers \((0),(1),(2)\), and so on. The same primitive root has always been chosen as base \(r\) as in Table I. Thus the rules of the preceding article give the period of any fraction whose denominator is in this table, after the numerator's index has been calculated from Table I. For denominators so small, however, we can complete the task just as easily without Table I. By ordinary division we calculate as many initial digits of the required mantissa as Article \(313\) requires to distinguish it from all others with the same denominator, no more than \(2\) for Table III. We then examine all the periods corresponding to the given denominator until we reach those initial digits, which will undoubtedly mark the beginning of the period. But we must point out that these digits may also be separated: the first, or several of them, may form the end of a period, and the remaining digit or digits its beginning.

\textit{Example.} We seek the period of \(\frac{12}{19}\). Here, for the modulus \(19\), Table I
\SourceLeaf{391}
gives \(\operatorname{ind.}12=2\operatorname{ind.}2+\operatorname{ind.}3=39\equiv3\; (\mathrm{mod}.\ 18)\), by Article \(57\). In this case there is just one period, corresponding to the numerator \(1\), so we must move its first three digits to the end. This gives the required period \(631578947368421052\). — The beginning of the period could have been found just as easily from its first two digits \(63\).

If we want the period of \(\frac{45}{53}\), then, for the modulus \(53\), \(\operatorname{ind.}45=2\operatorname{ind.}3 +\operatorname{ind.}5=49\). Here the number of periods is \(4=f\), and \(49=12f+1\), so, from the period marked \((1)\), the first \(12\) digits must be placed after the last digit. The required period is therefore \(8490566037735\). In this case the initial digits \(84\) are separated in the table.

We further observe that Table III can also give the number corresponding to a given index for a given modulus, listed there under the heading of denominator, as we promised in Article \(59\). By the preceding results, we can clearly find the period of the fraction whose numerator, though unknown, corresponds to the given index. It is enough to take as many initial digits of this period as the denominator has digits. From them, Article \(313\) gives the numerator, or the required number corresponding to the given index.

\Needspace{7\baselineskip}
\Article{317}
The preceding results give, without calculation, as many digits as we want of the mantissa of any fraction whose denominator is a prime or a power of a prime within the table's limits. But the investigations at the beginning of this section extend the table's range much further, to all fractions whose denominators are products of primes or powers of primes within those limits. For such a fraction can be decomposed into others whose denominators are these factors, and these can be converted into decimals to any number of digits; it remains only to add them. We hardly need to point out that the last digit of the resulting sum may be less than the correct one. But the shortfall clearly cannot reach as many units as there are individual fractions being added. It is therefore useful to calculate these to a few more digits than we want the given fraction to be correct to. As an example, consider the fraction \(\frac{6099380351}{1271808720}=F\)\HistoricalNote{*)}{This fraction is one of those that approximate the square root of \(23\) as closely as possible; its excess is less than seven units in the twentieth decimal place.}, whose denominator is the product of the numbers
\SourceLeaf{392}
\(16,\allowbreak 9,\allowbreak 5,\allowbreak 49,\allowbreak 13,\allowbreak 47,\allowbreak 59\). By the rules above we find \(F=1+\frac{11}{16}+\frac49+\frac45+\frac{22}{49}+\frac5{13}+\frac7{47}+\frac{52}{59}\). These individual fractions are converted into decimals as follows:
\[\renewcommand{\arraystretch}{1.25}\begin{array}{r@{\;}c@{\;}l@{\quad}l@{\quad}l}
1&{}={}&1&&\\
\frac{11}{16}&{}={}&0,6875&&\\
\frac45&{}={}&0,8&&\\
\frac49&{}={}&0,4444444444&4444444444&44\\
\frac{22}{49}&{}={}&0,4489795918&3673469387&75\\
\frac5{13}&{}={}&0,3846153846&1538461538&46\\
\frac7{47}&{}={}&0,1489361702&1276595744&68\\
\frac{52}{59}&{}={}&0,8813559322&0338983050&84\\\hline
F&{}={}&4,7958315233&1271954166&17
\end{array}\]
The shortfall of this sum from the correct value is certainly less than five units in the last, twenty-second, digit, so the first twenty-one digits cannot be changed by it. Extending the calculation to more digits replaces the last two digits \(17\) by \(1893936\ldots\). — Anyone will see, even without our pointing it out, that this method of converting common fractions into decimals is particularly suited to the case where many decimal digits are wanted. When a few suffice, ordinary division or logarithms can usually be used just as conveniently.

\Needspace{28\baselineskip}
\Article{318}
We have thus reduced the conversion of fractions whose denominators are composed of several distinct primes to the case where the denominator is a prime or a power of a prime. So we will add only a few observations about their mantissas. If the denominator contains neither the factor \(2\) nor \(5\), the mantissa here too consists of periods. For in this case as well the series \(10,\allowbreak 100,\allowbreak 1000\) eventually reaches a term congruent to unity modulo the denominator. The exponent of this term, which can easily be determined by Article \(92\), also gives the length of the period, independently of the numerator, provided the latter is coprime to the denominator. — But if the denominator has the form \(2^\alpha5^\beta N\), where \(N\) is coprime to \(10\), and \(\alpha\) and \(\beta\) are numbers at least one of which is not \(0\), the mantissa begins to consist of periods after the first \(\alpha\) or \(\beta\) digits, according as \(\alpha\) or \(\beta\) is greater. These periods agree in length with those of fractions with denominator \(N\).
\SourceLeaf{393}
This follows very easily from the decomposition of that fraction into two others with denominators \(2^\alpha5^\beta\) and \(N\), the former of which terminates after the first \(\alpha\) or \(\beta\) digits. — We could add many other observations on this subject, especially on devices for constructing a table like III as quickly as possible. For brevity we leave them out here all the more readily because much on this subject has already been given both by the distinguished Robertson in the work cited and by the distinguished Bernoulli (\textit{Nouv. Mém. de l'Ac. de Berlin} \(1771\ p.\ 273\)).
